\section{Một vector cho cả câu}
\label{sec:ch05-mot-vector}

\index{vector ngữ cảnh!định nghĩa}%
\tn{Mạng hồi quy}{recurrent neural network} của chương~\ref{ch:ky-hieu-va-bptt}
ánh xạ chuỗi sang chuỗi được, nhưng chỉ khi hai chuỗi khớp nhau từng bước: mỗi
\(\vect{x}_t\) cho một \(\vect{y}_t\). Dịch không như thế.
\tn{Câu nguồn}{source sentence} dài \(T\) token,
\tn{câu đích}{target sentence} dài \(T'\), hai số ấy khác nhau, và bạn không
biết \(T'\) trước khi dịch xong.

Bài giải bằng cách tách bài toán làm hai. Viết \(\vect{v}\) cho
\tn{vector ngữ cảnh}{context vector}, tức trạng thái cuối của encoder sau khi
nó đọc hết câu nguồn. Khi ấy

\begin{equation}
  p(\vect{y}_1, \dots, \vect{y}_{T'} \mid \vect{x}_1, \dots, \vect{x}_T)
  = \prod_{t=1}^{T'} p(\vect{y}_t \mid \vect{v}, \vect{y}_1, \dots, \vect{y}_{t-1}),
  \label{eq:ch05-phan-ra}
\end{equation}

đúng phương trình (1) của bài. Đọc nó chậm, vì cả chương nằm trong đó. Vế phải
không còn nhắc tới \(\vect{x}\) nữa. Sau khi encoder chạy xong, câu nguồn biến
mất khỏi bài toán; thứ duy nhất decoder còn thấy là \(\vect{v}\) và những gì
chính nó đã viết ra.

\index{vector ngữ cảnh!bề rộng cố định}%
Và \(\vect{v}\) có bề rộng cố định. Câu nguồn năm token hay năm mươi token thì
\(\vect{v}\) vẫn là chừng ấy số thực. Bài nói thẳng con số của mô hình họ: bốn
tầng, mỗi tầng một nghìn cell, tức tám nghìn số thực cho một câu. Nguyên văn:
\enquote{the deep LSTM uses 8000 real numbers to represent a sentence}.

\subsection*{Ký hiệu \(\langle\)EOS\(\rangle\) làm một việc cụ thể}

\index{vector ngữ cảnh!ký hiệu kết thúc}%
Mỗi câu kết thúc bằng một token đặc biệt, viết là \(\langle\)EOS\(\rangle\).
Không phải quy ước cho gọn. Bài nêu đúng lý do: ký hiệu ấy cho mô hình định
nghĩa được một phân phối trên các chuỗi thuộc mọi độ dài, nguyên văn
\enquote{enables the model to define a distribution over sequences of all
possible lengths}. Không có nó thì
\eqref{eq:ch05-phan-ra} là phân phối trên các chuỗi độ dài \(T'\) cố định, mà
\(T'\) thì mô hình phải tự quyết. Có nó thì độ dài trở thành một thứ mô hình
sinh ra, giống như sinh ra một từ.

\subsection*{Ba chỗ mô hình thật khác mô tả trên}

Mục 2 của bài liệt kê ba chỗ, và cả ba đều đáng đọc vì cả ba đều còn sống tới
hôm nay.

Thứ nhất, \emph{hai} LSTM chứ không phải một: một cho câu nguồn, một cho câu
đích, trọng số riêng biệt. Lý do bài đưa ra là nó tăng số tham số mà gần như
không tốn thêm tính toán, nguyên văn \enquote{at negligible computational
cost}, và làm cho việc huấn luyện nhiều cặp ngôn ngữ cùng lúc trở nên tự nhiên.
Trong repo companion, đó là hai module \code{encoder} và \code{decoder} độc lập
nhau.

Thứ hai, bốn tầng. Bài đo được mỗi tầng thêm vào giảm perplexity gần mười phần
trăm.

Thứ ba, đảo thứ tự từ của câu nguồn. Đây là chỗ tôi để dành hẳn
mục~\ref{sec:ch05-dao-cau-nguon}, vì nó là thứ duy nhất trong bài không đụng
tới kiến trúc mà vẫn đổi kết quả nhiều nhất.

\subsection*{\enquote{LSTM formulation from Graves} là bản nào}

\index{ô nhớ!bản Graves 2013}%
Bài viết đúng một câu về ô nhớ nó dùng: \enquote{We use the LSTM formulation
from Graves}~\autocite{graves2013generating}. Câu ấy đáng tra, vì
chương~\ref{ch:lstm} vừa dành cả chương để chỉ ra rằng \enquote{LSTM} không
phải một kiến trúc mà là vài kiến trúc.

Tra ra thì bản Graves dùng, mục 2.1 phương trình (7) tới (11), có
\tn{cổng quên}{forget gate} của Gers 2000, có peephole của Gers 2002, và dùng
tanh ở cả hai chỗ nén. Nghĩa là đúng hai hộp cầu nối mà mục~\ref{sec:ch04-cau-noi}
đã dựng, cộng lại. Graves ghi rõ các ma trận từ cell sang cổng là ma trận chéo.
Ông cũng nói thẳng chỗ mà mục~\ref{sec:ch04-cat-gradient} đã phân tích: bản gốc 1997
dùng một phép tính gradient xấp xỉ tự thiết kế, còn bài của ông tính gradient
đầy đủ bằng lan truyền ngược qua thời gian.

Repo companion cài bản có cổng quên và tanh, không có peephole, và một tầng
thay vì bốn. Hai chỗ bớt đi ấy là quyết định về kích thước để mọi thí nghiệm
chạy xong trên CPU trong ngân sách của sách, không phải bất đồng với bài; tôi
nói ra đây thay vì để listing ngầm bảo rằng bài đơn giản hơn thực tế.

\begin{minted}{python}
c = f * c_prev + i * torch.tanh(g)
h = o * torch.tanh(c)
\end{minted}

\index{vector ngữ cảnh!đường truyền duy nhất}%
Chỗ nối giữa hai mạng là hai hàm, và đọc chúng cạnh nhau thì thấy ngay
\(\vect{v}\) hẹp tới mức nào. Encoder trả về một cặp trạng thái; decoder nhận
đúng cặp ấy làm trạng thái khởi tạo rồi không bao giờ nhìn lại câu nguồn nữa:

\begin{minted}{python}
def encode(
    self, source: torch.Tensor
) -> tuple[torch.Tensor, torch.Tensor]:
    projected = self.encoder.project(self.source_embedding(source))
    state = self.encoder.initial_state(source.shape[1], projected.dtype)
    for t in range(source.shape[0]):
        live = (source[t] != PAD_INDEX).unsqueeze(-1).to(projected.dtype)
        c_next, h_next = self.encoder.step_projected(projected[t], state)
        state = (
            live * c_next + (1.0 - live) * state[0],
            live * h_next + (1.0 - live) * state[1],
        )
    return state
\end{minted}

\begin{minted}{python}
def decode_forced(
    self,
    context: tuple[torch.Tensor, torch.Tensor],
    target_in: torch.Tensor,
) -> torch.Tensor:
    projected = self.decoder.project(self.target_embedding(target_in))
    state = context
    outputs = []
    for t in range(target_in.shape[0]):
        state = self.decoder.step_projected(projected[t], state)
        outputs.append(state[1])
    return self.readout(torch.stack(outputs))
\end{minted}

Hai dòng \code{live} trong \code{encode} là chỗ dễ bỏ sót nhất trong cả module.
Chúng đóng băng trạng thái ở mọi vị trí mà input là token đệm, nên một câu ngắn
nằm chung batch với câu dài vẫn kết thúc ở trạng thái tại từ cuối của
\emph{chính nó}. Bỏ hai dòng ấy đi thì \(\vect{v}\) của một câu phụ thuộc vào
việc nó bị xếp chung batch với ai, hàm loss không hề kêu, và cái duy nhất bạn
thấy là một độ chính xác nhúc nhích khi đổi batch size. Repo có test riêng khoá
chuyện đó lại.

Hình~\ref{fig:ch05-encoder-decoder} vẽ toàn bộ kiến trúc.

\begin{figure}[htbp]
  \centering
  % Encoder-decoder của Sutskever 2014, vẽ theo hình 1 của bài.
% Câu nguồn đã đảo, nên encoder đọc "C B A" và token gần chỗ nối nhất là A,
% tức từ đầu tiên của câu gốc.
% Mono-safe: chỉ dùng nét liền, nét đứt, độ dày nét và mức xám.
\begin{tikzpicture}[
  >= Stealth,
  buoc/.style    = {draw, rectangle, minimum width = 8mm, minimum height = 6mm,
                    inner sep = 0pt},
  mui/.style     = {->, line width = 0.6pt, black!75},
  ngucanh/.style = {->, line width = 1.6pt, black},
  nho/.style     = {font = \scriptsize},
  ghichu/.style  = {font = \footnotesize, black!60},
]

  % --- encoder: ba bước, đọc câu nguồn đã đảo
  \foreach \i/\tok in {0/C, 1/B, 2/A} {
    \node[buoc] (e\i) at (\i*1.35, 0) {};
    \draw[mui] (\i*1.35, -1.0) -- (e\i);
    \node[nho, anchor = north] at (\i*1.35, -1.05) {\tok};
  }
  \draw[mui] (e0) -- (e1);
  \draw[mui] (e1) -- (e2);

  % --- chỗ nối: vector ngữ cảnh
  \draw[ngucanh] (e2) -- (4.75, 0) node[midway, above, nho] {\(\vect{v}\)};

  % --- decoder: bốn bước, sinh câu đích
  \foreach \i/\vao/\ra in {0/{\(\langle\)SOS\(\rangle\)}/X, 1/X/Y, 2/Y/Z,
                           3/Z/{\(\langle\)EOS\(\rangle\)}} {
    \pgfmathsetmacro\x{4.75 + \i*1.35}
    \node[buoc] (d\i) at (\x, 0) {};
    \draw[mui] (\x, -1.0) -- (d\i);
    \node[nho, anchor = north] at (\x, -1.05) {\vao};
    \draw[mui] (d\i) -- (\x, 1.0);
    \node[nho, anchor = south] at (\x, 1.05) {\ra};
  }
  \draw[mui] (d0) -- (d1);
  \draw[mui] (d1) -- (d2);
  \draw[mui] (d2) -- (d3);

  % --- Không vẽ mũi tên quay lui từ output về input bước sau. Hình 1 của bài
  % cũng không vẽ, và vì một lý do tốt: hai hàng nhãn đã nói điều đó rồi (ra
  % X Y Z, vào <sos> X Y Z), còn ba đường quay lui thì chỉ đi lọt được bằng
  % cách cắt qua chính các node nó nối.

  % --- khung và nhãn hai nửa
  \begin{scope}[on background layer]
    \draw[black!45, densely dotted, line width = 0.5pt]
      (-0.75, -0.75) rectangle (3.45, 0.75);
    \draw[black!45, densely dotted, line width = 0.5pt]
      (4.05, -0.75) rectangle (8.65, 0.75);
  \end{scope}
  % Cả hai nhãn khung đặt trên hàng token output, không đặt sát khung: ở cao độ
  % của khung thì chữ "decoder" rơi đúng vào giữa hai token Y và Z.
  \node[ghichu, anchor = south] at (1.35, 1.6) {encoder};
  \node[ghichu, anchor = south] at (6.35, 1.6) {decoder};

  \node[nho, anchor = north, black!60, align = center] at (1.35, -1.75)
    {câu nguồn, đã đảo\\ (gốc là A B C)};
  \node[nho, anchor = north, black!60] at (6.35, -1.75) {câu đích};

\end{tikzpicture}

  \caption{Encoder-decoder, vẽ theo hình 1 của bài. Encoder đọc câu nguồn đã
    đảo và không sinh ra gì. Decoder sinh từ, và hai hàng nhãn của nó đọc lệch
    nhau một bước vì mỗi từ nó sinh ra quay lại làm input của bước sau: hàng
    trên là X, Y, Z, \(\langle\)EOS\(\rangle\) và hàng dưới là
    \(\langle\)SOS\(\rangle\), X, Y, Z. Chỗ duy nhất nối hai nửa là mũi tên đậm
    ở giữa, mang \(\vect{v}\): bề rộng của nó cố định, nên toàn bộ câu nguồn
    phải đi lọt qua đó dù dài bao nhiêu. Chương~\ref{ch:dong-chinh} là chương
    nối thêm đường thứ hai qua chỗ thắt (bottleneck) này.}
  \label{fig:ch05-encoder-decoder}
\end{figure}
